(1) 连续两个波峰在声源参考钟上的发射时间差为 $\Delta t=1/f_0$。在这段时间内前方第一个波峰传播 $c/f_0$，而源前进 $u_s/f_0$，所以前方波长 $\lambda_{\rm front}=(c-u_s)/f_0$。静止观察者接收频率 $f=c/\lambda_{\rm front}=f_0c/(c-u_s)$。
(2) 若观察者沿声向运动的速度为 $u_o$，相对于波峰的接近速度为 $c-u_o$，每秒跨过的波峰数为 $f=(c-u_o)/\lambda_{\rm front}$，所以 $f=f_0(c-u_o)/(c-u_s)$。观察者逆着传播方向运动时 $u_o<0$，会提高接收频率。
(3) 声源远离观察者时后方相邻波峰间隔为 $\lambda_{\rm back}=(c+u_s)/f_0$，所以静止观察者测得 $f_{\rm away}=f_0c/(c+u_s)$。前方频率 $f_{\rm approach}=f_0c/(c-u_s)$ 较高，两者差别完全由波峰间距压缩与拉长造成。
(4) $u_s,u_o\to0$ 时各式都趋回 $f_0$。当 $u_s\to c^-$，前方线性模型预言波峰间距趋零，但实际有限振幅导致激波及非线性效应；超声速源有Mach锥，不应延用简单的亚声速周期波峰公式。
